\sgaStableTarget{sga6.e04.expose}\section*{第 IV 讲. 赋环拓扑斯的 Grothendieck 群}
\begin{center}
L. Illusie
\end{center}

\sgaStableTarget{sga6.e04.sec.1}\section*{1. Grothendieck 群的回顾与一般性质}

\sgaStableTarget{sga6.e04.para.1-1}\textbf{1.1.} 设 \(C\) 为三角范畴. 回顾一下（\hyperlink{sga6.e01.para.6-3}{I 6.3} 及 SGA 5 VIII 2）, 从 \(\Ob C\) 到 Abel 群 \(G\) 的映射 \(f\) 称为\emph{可加的（additive）}, 如果等式
\[
 f(E)=f(E')+f(E'')
\]
对每个特殊三角形
\[
 E'\longrightarrow E\longrightarrow E''\longrightarrow E'[1].
\]
都成立. 从 \(\Ob C\) 到 \(G\) 的可加映射构成一个 Abel 群, 它函子性地依赖于 \(G\). 所得函子由一个 Abel 群 \(k(C)\) 及一个泛可加映射
\[
 \cl_C:\Ob C\longrightarrow k(C),
\]
表示; 在不致混淆时, 将此映射简记为 \(\cl\).

群 \(k(C)\) 关于正合函子函子性地依赖于 \(C\). 若 \(C,C'\) 为三角范畴, 则从 \(C\) 到 \(C'\) 的两个同构正合函子诱导相同的同态 \(k(C)\to k(C')\); 特别地, 若 \(u:C\to C'\) 为正合范畴等价, 则 \(k(u):k(C)\to k(C')\) 为同构.

\sgaStableTarget{sga6.e04.lem.1-2}\textbf{引理 1.2.} 设 \(C\) 为三角范畴, \(L\in\Ob C\).
\begin{enumerate}[label=\alph*)]
\item 对任意 \(n\in\mathbb Z\), 有
\[
 \cl(L[n])=(-1)^n\cl(L).
\]
若 \(L',L''\) 为 \(C\) 的对象, 且 \(L\simeq L'\oplus L''\), 则
\[
 \cl(L)=\cl(L')+\cl(L'').
\]
\item 假设存在 \(n\in\mathbb Z\), 使得
\[
 \coprod_{i\geq0} L[2ni]
\]
可表. 则 \(\cl(L)=0\).
\end{enumerate}

\emph{证明.} 断言 a) 直接由定义得到（参见 SGA 5 VIII 2）. 对于 b), 只需注意到
\[
 \coprod_{i\geq0} L[2ni]\simeq
 L\oplus\left(\coprod_{i\geq0}L[2ni]\right)[2n],
\]
并应用 a).

\sgaStableTarget{sga6.e04.lem.1-3}\textbf{引理 1.3.} 设 \(C\) 为三角范畴, \(A\) 为 \(C\) 的厚子范畴 \(([\mathrm V]\ \mathrm{I}\ 2.1)\). 包含函子与商函子定义一个正合列
\[
 k(A)\longrightarrow k(C)\longrightarrow k(C/A)\longrightarrow 0.
\]

\emph{证明.} 见 SGA 5 VIII 3.1.

\sgaStableTarget{sga6.e04.lem.1-4}\textbf{引理 1.4.} 设 \(A\) 为可加范畴（分别为 Abel 范畴）, \(f\) 为 \(\Ob K^b(A)\) 上（分别为 \(\Ob D^b(A)\) 上）的可加函数（\hyperlink{sga6.e04.para.1-1}{1.1}）. 对 \(E\in\Ob K^b(A)\)（分别对 \(E\in\Ob D^b(A)\)）, 有
\[
 f(E)=\sum_i(-1)^i f(E^i),
\]
分别有
\[
 f(E)=\sum_i(-1)^i f(H^i(E)).
\]

\emph{证明.} 留给读者作为习题.

\sgaStableTarget{sga6.e04.para.1-5}\textbf{1.5.} 设 \(C\) 为 Abel 范畴, \(A\) 为 \(C\) 中对有限和稳定的满子范畴. 从 \(\Ob A\) 到 Abel 群 \(G\) 的映射 \(f\) 称为\emph{可加的}, 如果等式
\[
 \sum_i(-1)^i f(E^i)=0
\]
对 \(C\) 中每个分量属于 \(A\) 的有界无上同调复形 \(E\) 都成立. 当 \(A\) 对满态射的核稳定时（\hyperlink{sga6.e01.sec.1.2}{I 1.2}）, 这等价于说, 等式
\[
 f(L)=f(L')+f(L'')
\]
对 \(C\) 中每个各项属于 \(A\) 的正合列
\[
 0\longrightarrow L'\longrightarrow L\longrightarrow L''\longrightarrow0
\]
都成立. 从 \(\Ob A\) 到 \(G\) 的可加映射构成一个函子性地依赖于 \(G\) 的 Abel 群. 用 \(K^{b,\phi}(A)\) 表示 \(K^b(A)\) 中由无上同调复形构成的厚子范畴. 由 \hyperlink{sga6.e04.lem.1-3}{(1.3)} 和 \hyperlink{sga6.e04.lem.1-4}{(1.4)} 容易得到, 由 \(\cl_{K^b(A)/K^{b,\phi}(A)}\) 诱导的映射
\sgaStableTarget{sga6.e04.eq.1-5-1}\begin{equation}\tag{1.5.1}
 \cl_A:\Ob A\longrightarrow k\bigl(K^b(A)/K^{b,\phi}(A)\bigr)
\end{equation}
为泛可加映射, 即二元组
\[
 \left(k\bigl(K^b(A)/K^{b,\phi}(A)\bigr),\cl_A\right)
\]
表示将每个 Abel 群 \(G\) 对应到从 \(\Ob A\) 到 \(G\) 的可加映射所成 Abel 群的函子. 记
\sgaStableTarget{sga6.e04.eq.1-5-2}\begin{equation}\tag{1.5.2}
 k\bigl(K^b(A)/K^{b,\phi}(A)\bigr)=k(A),
\end{equation}
在需要避免环境 Abel 范畴的歧义时, 将 \(k(A)\) 写作 \(k(A,C)\).\footnote{在 \hyperlink{sga6.e01.para.6-3}{I 6.3} 中, 曾用记号 \(k(A)\) 表示群 \(k(K^b(A))\); 一般而言, 它与 \hyperlink{sga6.e04.eq.1-5-2}{(1.5.2)} 定义的群并不同构. 因此, 这里改变约定: 本讲中记号 \(k(A)\) 只表示由 \hyperlink{sga6.e04.eq.1-5-2}{(1.5.2)} 定义的群.}
特别地, 当 \(A=C\) 时, 有
\sgaStableTarget{sga6.e04.eq.1-5-3}\begin{equation}\tag{1.5.3}
 k(D^b(C))=k(C).
\end{equation}

\sgaStableTarget{sga6.e04.lem.1-6}\textbf{引理 1.6.} 设 \(C\) 为 Abel 范畴, \(A\) 为厚 Abel 子范畴, 即对核、余核和扩张稳定的满子范畴. 设 \(D_A(C)\) 为 \(D(C)\) 中由满足对每个 \(i\) 都有 \(H^i(E)\in\Ob A\) 的对象 \(E\) 构成的满子范畴. 则 \(D_A(C)\) 为 \(D(C)\) 的三角子范畴 \(([\mathrm V]\ \mathrm I\ 1.2.3)\), 且典范函子
\[
 D^b(A)\longrightarrow D_A^b(C)=D^b(C)\cap D_A(C)
\]
诱导一个同构
\[
 k(A)\xrightarrow{\sim} k(D_A^b(C)),
\]
其中 \(k(A)\) 由 \hyperlink{sga6.e04.eq.1-5-2}{(1.5.2)} 定义.

\emph{证明.} 第一个断言直接由上同调长正合列得到. 对于第二个断言, 考虑映射
\[
 f:\Ob D_A^b(C)\longrightarrow k(A)
\]
其定义为
\[
 f(E)=\sum_i(-1)^i\cl_A(H^i(E)).
\]
其中 \(\cl_A\) 为映射 \hyperlink{sga6.e04.eq.1-5-1}{(1.5.1)}. 由上同调长正合列, \(f\) 在 \hyperlink{sga6.e04.para.1-1}{(1.1)} 的意义下可加, 因而定义一个同态
\[
 \varphi:k(D_A^b(C))\longrightarrow k(A).
\]
由 \hyperlink{sga6.e04.lem.1-4}{(1.4)} 立即得到, \(\varphi\) 为典范同态 \(k(A)\to k(D_A^b(C))\) 的逆.

\sgaStableTarget{sga6.e04.rem.1-6-1}\textbf{注 1.6.1.} 在 \hyperlink{sga6.e04.lem.1-6}{(1.6)} 的情形下, 典范函子
\[
 D^b(A)\longrightarrow D_A^b(C)
\]
一般不是范畴等价. 取 \(S\) 为概形, \(C\) 为 \(\mathcal O_S\)-模范畴, \(A\) 为拟凝聚 \(\mathcal O_S\)-模范畴, 即得一个例子（\hyperlink{sga6.e02.para.0-3}{II 附录 I 0.3}）. 不过, 可参见判别法（\hyperlink{sga6.e02.lem.2-1-3}{II 2.1.3}）.

\sgaStableTarget{sga6.e04.def.1-7}\textbf{定义 1.7.} 设 \(C\) 为三角范畴（分别为 Abel 范畴）, \(C'\) 为 \(\Ob C\) 的子集. 称 \(C'\) 为\emph{正合子集（exact subset）}（参见 EGA III 3.1.1）, 如果对每个特殊三角形
\[
 E'\longrightarrow E\longrightarrow E''\longrightarrow E'[1],
\]
分别对每个正合列
\[
 0\longrightarrow E'\longrightarrow E\longrightarrow E''\longrightarrow0,
\]
只要 \(E',E,E''\) 中有两项属于 \(C'\), 第三项也属于 \(C'\). 若 \(\Ob C\) 是 \(\Ob C\) 中包含 \(C'\) 的最小正合子集, 则称 \(C'\) 在 \(\Ob C\) 中\emph{稠密（dense）}.

\sgaStableTarget{sga6.e04.lem.1-8}\textbf{引理 1.8.} 设 \(C\) 为三角范畴（分别为 Abel 范畴）, \(C'\) 为 \(\Ob C\) 的稠密子集. 则 \(E\in C'\) 时的元素 \(\cl_C(E)\) 生成 \(k(C)\).

\emph{证明.} 置 \(C'=C_0\), 并对 \(i\geq1\) 递归定义 \(C_i\): 它是 \(C\) 中那些出现在某个特殊三角形（分别为短正合列）中的对象所成的集合, 该三角形（分别为短正合列）有两项属于 \(C_{i-1}\). 各 \(C_i\) 构成 \(\Ob C\) 的递增滤过, 由假设它是穷尽的. 对 \(i\geq0\), 设 \(k(C)_i\) 为元素 \(\cl_C(E)\), \(E\in C_i\) 在 \(k(C)\) 中生成的子群. 从定义立即得到, 对每个 \(i\geq0\), 有 \(k(C)_i=k(C)_{i+1}\), 从而 \(k(C)_i=k(C)_0\). 若 \(x\in k(C)\), 容易看出可写成 \(x=\cl(E')-\cl(E'')\), 其中 \(E',E''\in\Ob C\)（在前一种情形下, 甚至可以取 \(E''=0\)）. 由于滤过 \((C_i)\) 穷尽, 存在 \(i\geq0\), 使 \(E'\) 与 \(E''\) 属于 \(C_i\); 因而 \(x\in k(C)_i=k(C)_0\), 即得所需结论.

\sgaStableTarget{sga6.e04.exercise.1-9}\textbf{习题 1.9.} 设 \(C\) 为 Abel 范畴, \(A\) 为对有限和、子对象及商稳定的满子范畴, 且 \(\Ob A\) 在 \(\Ob C\) 中稠密. 则由 \(\Ob A\) 上的可加函数 \(E\mapsto\cl_C(E)\)（参见 \hyperlink{sga6.e04.para.1-5}{(1.5)}）定义的同态
\[
 k(A,C)\longrightarrow k(C)
\]
为同构.

\emph{提示.} 置 \(A=A_0\), 并对 \(i\geq1\) 定义 \(A_i\) 为 \(C\) 的满子范畴, 其对象是 \(C\) 中 \(A_{i-1}\) 的两个对象的扩张. 证明 \(A_i\) 满足与 \(A\) 相同的假设, 再证明典范同态
\[
 \varinjlim_i k(A_i,C)\longrightarrow k(C)
\]
为同构. 所需结论随即得到.

\sgaStableTarget{sga6.e04.sec.2}\section*{2. 赋环拓扑斯的函子 \(K_{\bullet}\) 与 \(K^{\bullet}\)}

\sgaStableTarget{sga6.e04.para.2-1}\textbf{2.1.} 设 \((X,A)\) 为赋环拓扑斯. 回顾以下记号:
\[
 D^b(A_X)_{\coh}
\]
为上同调有界的左 \(A\)-模伪凝聚复形范畴（\hyperlink{sga6.e01.ex.2-15}{I 2.15}）;
\[
 D(A_X)_{\underline{\parf}}
\]
为有限完美振幅的左 \(A\)-模复形范畴（\hyperlink{sga6.e01.rem.5-8-4}{I 5.8.4}）;
\[
 \Loclib(A_X)
\]
为局部是有限自由模的直和因子的左 \(A\)-模范畴（\hyperlink{sga6.e02.para.2-0}{II 2.0}）;
\[
 \Coh(A_X)
\]
为相对于有限型自由 \(A\)-模范畴的凝聚左 \(A\)-模范畴（\hyperlink{sga6.e01.def.3-1}{I 3.1}）.

后两个范畴是左 \(A\)-模范畴 \(\Mod(A_X)\) 的满子范畴, 前两个范畴则是以下范畴的满三角子范畴:
\[
 D(A_X)=D(\Mod(A_X)).
\]

\sgaStableTarget{sga6.e04.def.2-2}\textbf{定义 2.2.} 设 \((X,A)\) 为赋环拓扑斯. 除非另有说明, 记
\[
 K_{\bullet}(A_X)=k(D^b(A_X)_{\coh}),\qquad
 K^{\bullet}(A_X)=k(D(A_X)_{\underline{\parf}}),
\]
其中 \(k(\ )\) 由 \hyperlink{sga6.e04.para.1-1}{(1.1)} 定义, 并记
\[
 K_{\bullet}(A_X)_{\naif}=k(\Coh(A_X)),\qquad
 K^{\bullet}(A_X)_{\naif}=k(\Loclib(A_X)),
\]
其中 \(k(\ )\) 由 \hyperlink{sga6.e04.eq.1-5-2}{(1.5.2)} 定义, 环境 Abel 范畴为 \(\Mod(A_X)\).\footnote{在不致混淆时, 允许省略指标 \(A\). 有时也略去指标 ``naive'', 但届时会明确说明.} 称 \(K_{\bullet}(A_X)\)（分别称 \(K^{\bullet}(A_X)\)）为 \(X\) 上伪凝聚（分别为完美）\(A\)-模复形的 Grothendieck 群; 称 \(K_{\bullet}(A_X)_{\naif}\)（分别称 \(K^{\bullet}(A_X)_{\naif}\)）为凝聚（分别为局部自由）左 \(A\)-模的 Grothendieck 群.

\sgaStableTarget{sga6.e04.rem.2-3}\textbf{注 2.3.}
\begin{enumerate}[label=\alph*)]
\item 可以定义 \(K_{\bullet}(\ )\) 与 \(K^{\bullet}(\ )\) 的变体, 置
\[
 K_{\bullet}^{\loc}(A_X)=k(D^{\loc b}(A_X)_{\coh}),
\]
其中 ``loc \(b\)'' 意为“上同调局部有界”, 并置
\[
 K^{\bullet}_{\loc}(A_X)=k(D(A_X)_{\parf}).
\]
当 \(X\) 的终对象拟紧时, 这些群与 \hyperlink{sga6.e04.def.2-2}{(2.2)} 所定义的群一致; 但一般而言, 它们较不适于相交理论中的计算.
\item 除非 \(A\) 作为左 \(A\)-模是凝聚的, 否则群 \(K_{\bullet}(A_X)_{\naif}\) 几乎没有意义; 因为此时很可能只有极少的凝聚层存在.
\item 考虑“过大”的三角范畴的 \(k\) 没有什么意义. 例如, 由（\hyperlink{sga6.e04.lem.1-2}{1.2} b）立即推出关系式
\[
\begin{aligned}
 k(D(A_X))&=k(D^-(A_X))=k(D^+(A_X))\\
 &=k(D^b(A_X))=k(D^-(A_X)_{\coh})=0,
\end{aligned}
\]
其中最后一个等式是注意到, \(L\in\Ob D^-(A_X)_{\coh}\) 蕴含
\[
 \coprod_{i\geq0}L[2i]\in\Ob D^-(A_X)_{\coh}.
\]
\end{enumerate}

\sgaStableTarget{sga6.e04.para.2-4}\textbf{2.4. 凝聚环层的情形.} 设 \((X,A)\) 为赋环拓扑斯, 假设 \(A\) 作为左 \(A\)-模是凝聚的（\hyperlink{sga6.e01.def.3-1}{I 3.1} 和 \hyperlink{sga6.e01.ex.3-6}{3.6}）. 每个凝聚左 \(A\)-模都是伪凝聚复形（\hyperlink{sga6.e01.cor.3-5}{I 3.5}）, 因而有全忠实函子
\sgaStableTarget{sga6.e04.eq.2-4-1}\begin{equation}\tag{2.4.1}
 \Coh(A_X)\longrightarrow D^b(A_X)_{\coh}.
\end{equation}
此外, \(\Coh(A_X)\) 是 \(\Mod(A_X)\) 的厚子范畴（\hyperlink{sga6.e01.cor.3-4}{I 3.4}）, 且由（\hyperlink{sga6.e01.cor.3-5}{I 3.5}）, \(D^b(A_X)_{\coh}\) 为 \(D^b(A_X)\) 中由满足对每个 \(i\) 都有 \(H^i(L)\in\Ob\Coh(A_X)\) 的对象 \(L\) 构成的满子范畴. 因此, 由 \hyperlink{sga6.e04.lem.1-6}{(1.6)} 得到, 由 \hyperlink{sga6.e04.eq.2-4-1}{(2.4.1)} 定义的同态
\sgaStableTarget{sga6.e04.eq.2-4-2}\begin{equation}\tag{2.4.2}
 K_{\bullet}(A_X)_{\naif}\longrightarrow K_{\bullet}(A_X)
\end{equation}
为同构.

\sgaStableTarget{sga6.e04.para.2-5}\textbf{2.5. 同态 \(K^{\bullet}\to K_{\bullet}\).} 设 \((X,A)\) 为赋环拓扑斯. 包含函子
\[
 D(A_X)_{\underline{\parf}}\longrightarrow D^b(A_X)_{\coh}
\]
定义一个同态
\sgaStableTarget{sga6.e04.eq.2-5-1}\begin{equation}\tag{2.5.1}
 K^{\bullet}(A_X)\longrightarrow K_{\bullet}(A_X).
\end{equation}
它一般既不单也不满. 不过, 假设满足下列条件:
\begin{enumerate}[label=(\roman*)]
\item \(X\) 的终对象拟紧;
\item \(X\) 有足够多的点, 且对 \(X\) 的每个点 \(x\), 每个左 \(A_x\)-模均有有限 Tor 维数.
\end{enumerate}
则由（\hyperlink{sga6.e01.cor.5-10}{I 5.10}）, 有
\[
 D(A_X)_{\underline{\parf}}=D^b(A_X)_{\coh},
\]
从而 \hyperlink{sga6.e04.eq.2-5-1}{(2.5.1)} 为同构. 例如, 赋以正则局部环的紧空间满足假设 (i) 和 (ii).

\sgaStableTarget{sga6.e04.para.2-6}\textbf{2.6. 非交换情形下 \(K^{\bullet}\) 与 \(K^{\bullet}\) 之间的配对.} 设
\[
 f:(X',A')\longrightarrow (X,A)
\]
为赋环拓扑斯的态射, 用 \(\Mod(A',X',A)\) 表示兼为左 \(A'\)-模和右 \(f^{-1}(A)\)-模的双模范畴. 设
\[
 \Loclib_{A'}(A',X',A)
\]
为由作为 \(A'\) 上的对象属于 \(\Loclib(A',X')\) 的对象构成的满子范畴. 仿照 \hyperlink{sga6.e04.def.2-2}{(2.2)}, 置
\[
 K^{\bullet}_{\naif}(A',X',A)_{A'}=k(\Loclib_{A'}(A',X',A)).
\]
张量积
\[
 \otimes_A:\Mod(A',X',A)\times\Mod(A_X)\longrightarrow \Mod(A',X')
\]
给出对应
\[
 (E',E)\longmapsto E'\otimes_{f^{-1}(A)}f^{-1}(E),
\]
并诱导函子
\sgaStableTarget{sga6.e04.eq.2-6-1}\begin{equation}\tag{2.6.1}
 \Loclib_{A'}(A',X',A)\times\Loclib(A_X)\longrightarrow\Loclib(A',X').
\end{equation}
由 \hyperlink{sga6.e04.eq.2-6-1}{(2.6.1)} 得到的映射 \((E',E)\mapsto\cl(E'\otimes_A E)\) 在 \hyperlink{sga6.e04.para.1-5}{(1.5)} 的意义下为双可加的, 因而定义同态
\sgaStableTarget{sga6.e04.eq.2-6-2}\begin{equation}\tag{2.6.2}
 K^{\bullet}_{\naif}(A',X',A)_{A'}\otimes K^{\bullet}_{\naif}(A_X)\longrightarrow K^{\bullet}_{\naif}(A',X').
\end{equation}

设 \(D(A',X',A)\) 为 \(\Mod(A',X',A)\) 的导出范畴, 设
\[
 D(A',X',A)_{A'\text{-}\underline{\parf}}
\]
为由相对于 \(A'\), 即作为 \(D(A',X')\) 的对象, 具有有限完美振幅的复形构成的满子范畴. 如 \hyperlink{sga6.e04.def.2-2}{(2.2)} 中那样, 置
\[
 K^{\bullet}(A',X',A)_{A'}=k(D(A',X',A)_{A'\text{-}\underline{\parf}}).
\]
包含函子定义显然的同态
\sgaStableTarget{sga6.e04.eq.2-6-3-1}\sgaStableTarget{sga6.e04.eq.2-6-3-2}\sgaStableTarget{sga6.e04.eq.2-6-3-3}\begin{align}
 K^{\bullet}_{\naif}(A',X',A)_{A'}&\longrightarrow K^{\bullet}(A',X',A)_{A'}, \tag{2.6.3.1}\\
 K^{\bullet}_{\naif}(A_X)&\longrightarrow K^{\bullet}(A_X), \tag{2.6.3.2}\\
 K^{\bullet}_{\naif}(A',X')&\longrightarrow K^{\bullet}(A',X'). \tag{2.6.3.3}
\end{align}
另一方面, 已知（\hyperlink{sga6.e01.prop.4-18}{I 4.18}, \hyperlink{sga6.e01.cor.4-19}{4.19}）, 导出张量函子
\[
 \otimes_A^{\mathbf L}:D^-(A',X',A)\times D^-(A_X)\longrightarrow D^-(A',X')
\]
诱导一个正合函子
\sgaStableTarget{sga6.e04.eq.2-6-4}\begin{equation}\tag{2.6.4}
 D(A',X',A)_{A'\text{-}\underline{\parf}}\times D(A_X)_{\underline{\parf}}
 \longrightarrow D(A',X')_{\underline{\parf}}.
\end{equation}
后者定义同态
\sgaStableTarget{sga6.e04.eq.2-6-5}\begin{equation}\tag{2.6.5}
 K^{\bullet}(A',X',A)_{A'}\otimes K^{\bullet}(A_X)\longrightarrow K^{\bullet}(A',X').
\end{equation}
由于 \hyperlink{sga6.e04.eq.2-6-1}{(2.6.1)} 由 \hyperlink{sga6.e04.eq.2-6-4}{(2.6.4)} 诱导, 同态 \hyperlink{sga6.e04.eq.2-6-2}{(2.6.2)} 和 \hyperlink{sga6.e04.eq.2-6-5}{(2.6.5)} 与同态 \hyperlink{sga6.e04.eq.2-6-3-1}{(2.6.3.1)}, \hyperlink{sga6.e04.eq.2-6-3-2}{(2.6.3.2)} 及 \hyperlink{sga6.e04.eq.2-6-3-3}{(2.6.3.3)} 相容; 换言之, 有交换图
\sgaStableTarget{sga6.e04.eq.2-6-6}\[
\begin{tikzcd}[column sep=tiny]
K^{\bullet}_{\naif}(A',X',A)_{A'}\otimes K^{\bullet}_{\naif}(A_X)\arrow[r]\arrow[d]
&K^{\bullet}_{\naif}(A',X')\arrow[d]\\
K^{\bullet}(A',X',A)_{A'}\otimes K^{\bullet}(A_X)\arrow[r]
&K^{\bullet}(A',X'),
\end{tikzcd}
\tag{2.6.6}
\]
其中上、下水平箭头分别为 \hyperlink{sga6.e04.eq.2-6-2}{(2.6.2)} 和 \hyperlink{sga6.e04.eq.2-6-5}{(2.6.5)}, 竖直箭头由 \hyperlink{sga6.e04.eq.2-6-3-1}{(2.6.3.1)}, \hyperlink{sga6.e04.eq.2-6-3-2}{(2.6.3.2)} 及 \hyperlink{sga6.e04.eq.2-6-3-3}{(2.6.3.3)} 定义.

\sgaStableTarget{sga6.e04.para.2-7}\textbf{2.7. \(K^{\bullet}\) 的反变性及环结构.} 下面详述配对 \hyperlink{sga6.e04.para.2-6}{(2.6)} 的一些重要特例, 虽然直接研究它们也是容易的.
\begin{enumerate}[label=\alph*)]
\item 设
\[
 f:(X',A')\longrightarrow (X,A)
\]
为赋环拓扑斯的态射. 函子
\[
 Lf^*:D^-(A_X)\longrightarrow D^-(A',X')
\]
将 \(D(A_X)_{\underline{\parf}}\) 映入 \(D(A',X')_{\underline{\parf}}\)（\hyperlink{sga6.e01.cor.4-19-1}{I 4.19.1}）, 同样将 \(\Loclib(A_X)\) 映入 \(\Loclib(A',X')\). 因此它定义一个同态
\sgaStableTarget{sga6.e04.eq.2-7-1}\begin{equation}\tag{2.7.1}
 f^*:K^{\bullet}(A_X)\longrightarrow K^{\bullet}(A',X')
\end{equation}
以及相应的同态
\sgaStableTarget{sga6.e04.eq.2-7-2}\begin{equation}\tag{2.7.2}
 f^*:K^{\bullet}(A_X)_{\naif}\longrightarrow K^{\bullet}(A',X')_{\naif}.
\end{equation}
这也可解释为在 \hyperlink{sga6.e04.eq.2-6-5}{(2.6.5)}（分别在 \hyperlink{sga6.e04.eq.2-6-2}{(2.6.2)}）中乘以 \(\cl(A')\).\footnote{译校注: 原书第 283 页此处将 (2.6.2) 与 (2.6.5) 的顺序互换了; 前者对应朴素群, 后者对应完美复形的群.} 有如下交换图, 它是 \hyperlink{sga6.e04.eq.2-6-6}{(2.6.6)} 的一个特例:
\[
\begin{tikzcd}
K^{\bullet}(A_X)_{\naif}\arrow[r,"f^*"]\arrow[d]
&K^{\bullet}(A',X')_{\naif}\arrow[d]\\
K^{\bullet}(A_X)\arrow[r,"f^*"]
&K^{\bullet}(A',X'),
\end{tikzcd}
\]
其中左、右竖直箭头分别为 \hyperlink{sga6.e04.eq.2-6-3-2}{(2.6.3.2)} 和 \hyperlink{sga6.e04.eq.2-6-3-3}{(2.6.3.3)}.

设 \(g:(X,A)\to(X'',A'')\) 为赋环拓扑斯的态射. 逆像函子的传递性蕴含, \hyperlink{sga6.e04.eq.2-7-1}{(2.7.1)} 型（分别为 \hyperlink{sga6.e04.eq.2-7-2}{(2.7.2)} 型）同态满足
\sgaStableTarget{sga6.e04.eq.2-7-3}\begin{equation}\tag{2.7.3}
 (gf)^*=f^*g^*.
\end{equation}
因而 \(K^{\bullet}(A_X)\)（分别为 \(K^{\bullet}(A_X)_{\naif}\)）为 \((X,A)\) 的反变函子.

\item 设 \((X,\mathcal O_X)\) 为赋交换环拓扑斯. 由于张量积的结合性与交换性, 取 \(A=\mathcal O_X\), \((X',A')=(X,A)\) 及 \(f=\id\) 后, \hyperlink{sga6.e04.eq.2-6-5}{(2.6.5)}（分别为 \hyperlink{sga6.e04.eq.2-6-2}{(2.6.2)}）赋予 \(K^{\bullet}(X)\)（分别为 \(K^{\bullet}(X)_{\naif}\)）带幺交换环结构, 其单位元为
\[
 1=\cl(\mathcal O_X).
\]
典范同态
\[
 K^{\bullet}(X)_{\naif}\longrightarrow K^{\bullet}(X)
\]
为带幺环的同态.

若 \(f:(X,\mathcal O_X)\to(Y,\mathcal O_Y)\) 为赋交换环拓扑斯的态射, 则张量积与逆像的交换性表明,
\[
 f^*:K^{\bullet}(Y)\longrightarrow K^{\bullet}(X),
\]
分别有
\[
 f^*:K^{\bullet}(Y)_{\naif}\longrightarrow K^{\bullet}(X)_{\naif},
\]
为带幺环的同态.
\end{enumerate}

\sgaStableTarget{sga6.e04.para.2-8}\textbf{2.8. \(K^{\bullet}\) 到 \(H^0(\ ,\mathbb Z)\) 的增广.} 设 \((X,\mathcal O_X)\) 为局部赋环拓扑斯（\hyperlink{sga6.e01.lem.2-15-1}{I 2.15.1}）. 回顾一下, 已经定义了（\hyperlink{sga6.e01.def.6-11}{I 6.11} 至 \hyperlink{sga6.e01.prop.6-14}{6.14}）带幺环同态
\[
 i:H^0(X,\mathbb Z)\longrightarrow K^{\bullet}(X)_{\naif},
\]
从而赋予 \(K^{\bullet}(X)_{\naif}\) 以 \(H^0(X,\mathbb Z)\) 上的带幺代数结构, 再经典范同态 \(K^{\bullet}(X)_{\naif}\to K^{\bullet}(X)\) 赋予 \(K^{\bullet}(X)\) 同样的结构. 此外, 秩给出 \(K^{\bullet}(X)_{\naif}\)（分别为 \(K^{\bullet}(X)\)）到 \(H^0(X,\mathbb Z)\) 的增广, 且典范同态 \(K^{\bullet}(X)_{\naif}\to K^{\bullet}(X)\) 为增广代数同态.

若 \(f:(X,\mathcal O_X)\to(Y,\mathcal O_Y)\) 为局部赋环拓扑斯的态射, 则同态
\[
 f^*:K^{\bullet}(Y)\longrightarrow K^{\bullet}(X),\qquad
 f^*:K^{\bullet}(Y)_{\naif}\longrightarrow K^{\bullet}(X)_{\naif}
\]
与上述增广代数结构相容（\hyperlink{sga6.e01.prop.6-12}{I 6.12}, \hyperlink{sga6.e01.para.6-13}{6.13}）; 这由下图的交换性表示, 其中 \(\Rg\) 表示秩:
\[
\begin{tikzcd}[column sep=tiny]
H^0(Y,\mathbb Z)\arrow[r,"i"]\arrow[d]
&K^{\bullet}(Y)_{\naif}\arrow[r]\arrow[d]
&K^{\bullet}(Y)\arrow[r,"\Rg"]\arrow[d]
&H^0(Y,\mathbb Z)\arrow[d]\\
H^0(X,\mathbb Z)\arrow[r,"i"]
&K^{\bullet}(X)_{\naif}\arrow[r]
&K^{\bullet}(X)\arrow[r,"\Rg"]
&H^0(X,\mathbb Z).
\end{tikzcd}
\]
\textsuperscript{*}在讨论班后文（\hyperlink{sga6.e05.ex.3-9-1}{V 3.9.1} 及 VI）将会看到, \(K^{\bullet}(X)_{\naif}\) 具有比到 \(H^0(X,\mathbb Z)\) 增广的代数更丰富的结构, 即到视作二项式环的 \(H^0(X,\mathbb Z)\)（\hyperlink{sga6.e05.ex.2-8}{V 2.8}）增广的 \(\lambda\)-代数结构; 此结构与逆像同态相容. Deligne [5]\footnote{译校注: 原书第 285 页此引文作 [5], 但本讲文献表中 Deligne 的条目为 [1].} 最近证明, 至少当 \(X\) 的特征为 \(p\geq0\) 时, \(K^{\bullet}(X)\) 具有类似结构, 且典范同态
\[
 K^{\bullet}(X)_{\naif}\longrightarrow K^{\bullet}(X)
\]
为增广 \(\lambda\)-代数同态, 从而肯定地回答了 Grothendieck 的一个猜想（XIV 1）.\textsuperscript{*}

\sgaStableTarget{sga6.e04.para.2-9}\textbf{2.9. \(K^{\bullet}\) 与 \(K^{\bullet}_{\naif}\) 的比较.} 设 \((X,A)\) 为赋环拓扑斯. 一般没有理由认为典范同态
\[
 K^{\bullet}(X)_{\naif}\longrightarrow K^{\bullet}(X)
\]
是单射或满射. 不过, \hyperlink{sga6.e02.expose}{第 II 讲}的“整体解消存在性判别法”表明, 在以下情形它是同构:
\begin{enumerate}[label=(\roman*)]
\item \(X\) 为分离或 Noether 除子性概形的 Zariski 拓扑斯, 且 \(A=\mathcal O_X\)（\hyperlink{sga6.e02.prop.2-2-9}{II 2.2.9}）;
\item \(X\) 为紧空间定义的拓扑斯, 且 \(A\) 为软层（\hyperlink{sga6.e02.prop.2-3-2}{II 2.3.2}）;
\item \(X\) 为终对象拟紧的 Stein 拓扑斯, 且 \(A=\mathcal O_X\)（\hyperlink{sga6.e02.prop.2-4-5}{II 2.4.5}）.
\end{enumerate}

\sgaStableTarget{sga6.e04.para.2-10}\textbf{2.10. \(K^{\bullet}\) 与 \(K_{\bullet}\) 之间的配对.} 设 \((X,\mathcal O_X)\) 为赋交换环拓扑斯. 已知（\hyperlink{sga6.e01.rem.5-8-4}{I 5.8.4} b）有配对
\[
 \otimes_{\mathcal O_X}^{\mathbf L}:D(X)_{\underline{\parf}}\times D^b(X)_{\coh}
 \longrightarrow D^b(X)_{\coh}.
\]
它定义同态
\sgaStableTarget{sga6.e04.eq.2-10-1}\begin{equation}\tag{2.10.1}
 K^{\bullet}(X)\otimes K_{\bullet}(X)\longrightarrow K_{\bullet}(X).
\end{equation}
由于张量积的结合性, \hyperlink{sga6.e04.eq.2-10-1}{(2.10.1)} 赋予 \(K_{\bullet}(X)\) 以 \(K^{\bullet}(X)\)-模结构. 更一般地, 若
\[
 f:(X,\mathcal O_X)\longrightarrow (Y,\mathcal O_Y)
\]
为赋交换环拓扑斯的态射, 则可经同态
\[
 f^*:K^{\bullet}(Y)\longrightarrow K^{\bullet}(X)
\]
将 \(K_{\bullet}(X)\) 视为 \(K^{\bullet}(Y)\)-模.

\sgaStableTarget{sga6.e04.para.2-11}\textbf{2.11. \(K_{\bullet}\) 的协变性与投影公式.} 设 \((X,A)\) 为赋环拓扑斯. 与 \(K^{\bullet}(A_X)\) 相比, 群 \(K_{\bullet}(A_X)\) 关于 \((X,A)\) 没有自然的变性. 然而, 在代数几何与解析几何中, \(K_{\bullet}(X)\) 对固有态射倾向于关于 \(X\) 协变. 下面予以精确说明.

\sgaStableTarget{sga6.e04.para.2-11-1}\textbf{2.11.1.} 设 \(f:X\to Y\) 为拟紧概形的固有伪凝聚态射（\hyperlink{sga6.e03.def.1-2}{III 1.2}）. 则只要 \(f\) 满足有限性猜想（\hyperlink{sga6.e03.conj.2-1}{III 2.1}）, 特别地（\hyperlink{sga6.e03.thm.2-2}{III 2.2}）, 若 \(Y\) 为 Noether 概形（此时 \(f\) 的伪凝聚性假设自动成立（\hyperlink{sga6.e03.ex.1-3}{III 1.3}））或 \(f\) 为射影态射, 由（\hyperlink{sga6.e03.schol.2-6}{III 2.6}）, 函子 \(Rf_*\) 诱导函子
\[
 Rf_*:D^b(X)_{\coh}\longrightarrow D^b(Y)_{\coh}.
\]
因此定义同态
\sgaStableTarget{sga6.e04.eq.2-11-1-1}\begin{equation}\tag{2.11.1.1}
 f_*:K_{\bullet}(X)\longrightarrow K_{\bullet}(Y).
\end{equation}
此同态为 \(K^{\bullet}(Y)\)-线性的（参见 \hyperlink{sga6.e04.para.2-10}{(2.10)}）; 换言之, 对 \(y\in K^{\bullet}(Y)\) 和 \(x\in K_{\bullet}(X)\), 有\emph{投影公式（projection formula）}
\sgaStableTarget{sga6.e04.eq.2-11-1-2}\begin{equation}\tag{2.11.1.2}
 f_*(yx)=y f_*(x).
\end{equation}
事实上, 若 \(x=\cl(E)\), \(E\in\Ob D^b(X)_{\coh}\), 且 \(y=\cl(M)\), \(M\in\Ob D(Y)_{\underline{\parf}}\), 则由（\hyperlink{sga6.e03.prop.3-7}{III 3.7}）有
\[
 Rf_*(E)\otimes_Y^{\mathbf L}M\simeq
 Rf_*(E\otimes_Y^{\mathbf L}M).
\]
最后, 若 \(g:Y\to Z\) 为满足与 \(f\) 相同假设的固有态射, 且有限性猜想（\hyperlink{sga6.e03.conj.2-1}{III 2.1}）对它成立, 则正像函子的传递性给出同态 \hyperlink{sga6.e04.eq.2-11-1-1}{(2.11.1.1)} 的公式
\sgaStableTarget{sga6.e04.eq.2-11-1-3}\begin{equation}\tag{2.11.1.3}
 g_*f_*=(gf)_*.
\end{equation}

\sgaStableTarget{sga6.e04.para.2-11-2}\textbf{2.11.2.} 设 \(f:X\to Y\) 为有限维解析空间的固有态射. 则由 Grauert 的有限性定理 [2], \(Rf_*\) 将 \(D^b(X)_{\coh}\) 映入 \(D^b(Y)_{\coh}\), 因而定义同态
\sgaStableTarget{sga6.e04.eq.2-11-2-1}\begin{equation}\tag{2.11.2.1}
 f_*:K_{\bullet}(X)\longrightarrow K_{\bullet}(Y).
\end{equation}
与代数情形的同态 \hyperlink{sga6.e04.eq.2-11-1-1}{(2.11.1.1)} 一样, 此同态为 \(K^{\bullet}(Y)\)-线性的; 等价地, 公式 \hyperlink{sga6.e04.eq.2-11-1-2}{(2.11.1.2)} 成立, 这由导出范畴层次的投影公式（\hyperlink{sga6.e03.rem.5-8}{III 5.9} b）得到. 此外, 还有传递公式 \hyperlink{sga6.e04.eq.2-11-1-3}{(2.11.1.3)}.

\sgaStableTarget{sga6.e04.para.2-11-3}\textbf{2.11.3.} 离开代数或解析的框架后, \(K_{\bullet}\) 对固有态射的这种协变性就不复存在. 例如, 设
\[
 f:X\to Y
\]
为赋以复值连续函数层的拓扑空间的固有态射. 则若 \(E\in\Ob D^b(X)_{\coh}\), 一般 \(Rf_*(E)\) 未必属于 \(D^b(Y)_{\coh}\), 即使取 \(E=\mathcal O_X\) 并令 \(Y\) 为一点, 这也已很明显. 实际上, \(K_{\bullet}\) 的\hyperlink{sga6.e04.def.2-2}{定义 (2.2)} 特别适合代数几何与解析几何的需要; 对拓扑空间, \(K_{\bullet}\) 的良好定义尚待发现. 此外, 看来必须离开 \(\mathcal O_X\)-线性复形的框架, 对拓扑空间或混合流形（II 附录 II）的态射 \(f:X\to Y\), 定义 \(X\) 上适当的一类拓扑向量层复形, 使得当 \(f\) 是混合流形的结构态射时, 这一类包含 \(X\) 上相对于 \(Y\) 的椭圆复形. 有限性定理（II 附录 II）及 Grothendieck 在 [3] 中对 \(Y\) 为一点的情形证明的类似定理, 都提示了这一方向.

\sgaStableTarget{sga6.e04.para.2-12}\textbf{2.12. \(K^{\bullet}\) 与 \(K_{\bullet}\) 的非常规变性.} 设
\[
 f:(X,A)\longrightarrow(Y,B)
\]
为赋环拓扑斯的有限 Tor 维数态射（\hyperlink{sga6.e03.def.3-1}{III 3.1}）. 函子
\[
 Lf^*:D^-(B_Y)\longrightarrow D^-(A_X)
\]
将 \(D^b(B_Y)\) 映入 \(D^b(A_X)\), 因而由（\hyperlink{sga6.e01.cor.2-16-1}{I 2.16.1}）, 将 \(D^b(B_Y)_{\coh}\) 映入 \(D^b(A_X)_{\coh}\), 于是定义同态
\sgaStableTarget{sga6.e04.eq.2-12-1}\begin{equation}\tag{2.12.1}
 f^*:K_{\bullet}(B_Y)\longrightarrow K_{\bullet}(A_X).
\end{equation}
若 \(A\) 与 \(B\) 交换, 则逆像与张量积的交换性表明, \hyperlink{sga6.e04.eq.2-12-1}{(2.12.1)} 与 \(K_{\bullet}(Y)\) 和 \(K_{\bullet}(X)\) 分别在 \(K^{\bullet}(Y)\) 和 \(K^{\bullet}(X)\) 上的模结构相容（\hyperlink{sga6.e04.para.2-10}{(2.10)}）. 换言之, 对 \(a\in K^{\bullet}(Y)\) 和 \(y\in K_{\bullet}(Y)\), 有
\sgaStableTarget{sga6.e04.eq.2-12-2}\begin{equation}\tag{2.12.2}
 f^*(ay)=a f^*(y).
\end{equation}
若 \(g:(Y,B)\to(Z,C)\) 为赋环拓扑斯的有限 Tor 维数态射, 则逆像函子的传递性蕴含, 同态 \hyperlink{sga6.e04.eq.2-12-1}{(2.12.1)} 满足
\sgaStableTarget{sga6.e04.eq.2-12-1-1}\begin{equation}\tag{2.12.1.1}
 f^*g^*=(gf)^*.
\end{equation}

现在假设处于 \hyperlink{sga6.e04.para.2-11-1}{(2.11.1)}（分别为 \hyperlink{sga6.e04.para.2-11-2}{(2.11.2)}）的情形, 且 \(f\) 有有限 Tor 维数. 则由（\hyperlink{sga6.e03.cor.4-8-1}{III 4.8.1} 及 \hyperlink{sga6.e03.rem.5-8}{5.9} b）, 函子 \(Rf_*\) 将 \(D(X)_{\underline{\parf}}\) 映入 \(D(Y)_{\underline{\parf}}\), 因而定义同态
\sgaStableTarget{sga6.e04.eq.2-12-3}\begin{equation}\tag{2.12.3}
 f_*:K^{\bullet}(X)\longrightarrow K^{\bullet}(Y).
\end{equation}
它一般不是环同态. 然而, 对 \(y\in K_{\bullet}(Y)\)（分别为 \(K^{\bullet}(Y)\)）及 \(x\in K^{\bullet}(X)\), 有
\sgaStableTarget{sga6.e04.eq.2-12-4}\begin{equation}\tag{2.12.4}
 f_*(f^*(y)x)=y f_*(x),
\end{equation}
其中 \(f^*\) 为同态 \hyperlink{sga6.e04.eq.2-12-1}{(2.12.1)}（分别为 \hyperlink{sga6.e04.eq.2-7-1}{(2.7.1)}）. 证明与 \hyperlink{sga6.e04.eq.2-11-1-2}{(2.11.1.2)} 类似. 若 \(g:Y\to Z\) 为满足与 \(f\) 相同假设的态射, 则同态 \hyperlink{sga6.e04.eq.2-12-3}{(2.12.3)} 满足传递公式
\sgaStableTarget{sga6.e04.eq.2-12-3-1}\begin{equation}\tag{2.12.3.1}
 g_*f_*=(gf)_*.
\end{equation}

无论在代数几何还是在解析几何中, 同态 \hyperlink{sga6.e04.eq.2-12-3}{(2.12.3)} 都是 Riemann--Roch 定理的基本组成部分之一; 解析情形的说明见\hyperlink{sga6.e00.sec.5}{第 0 讲, 5}.

假设 \(X\) 和 \(Y\) 为分离除子性概形（\hyperlink{sga6.e02.def.2-2-5}{II 2.2.5}）, 例如具有丰沛可逆层的概形. 则 \hyperlink{sga6.e04.para.2-9}{(2.9 (i))} 将 \(K^{\bullet}(X)\)（分别为 \(K^{\bullet}(Y)\)）与 \(K^{\bullet}(X)_{\naif}\)（分别为 \(K^{\bullet}(Y)_{\naif}\)）等同, 因而 \hyperlink{sga6.e04.eq.2-12-3}{(2.12.3)} 定义同态
\sgaStableTarget{sga6.e04.eq.2-12-3.occ2}\begin{equation}\tag{2.12.3'}
 f_*:K^{\bullet}(X)_{\naif}\longrightarrow K^{\bullet}(Y)_{\naif}.
\end{equation}
如\hyperlink{sga6.e00.sec.3.1}{第 0 讲, 3.1} 已说明的那样, 先验地看, （\hyperlink{sga6.e04.eq.2-12-3}{2.12.3}'）的定义远非显然. 实际上, 这正是本讨论班从\hyperlink{sga6.e01.expose}{第 I 讲}到\hyperlink{sga6.e04.expose}{第 IV 讲}的主要动机之一.

\sgaStableTarget{sga6.e04.sec.3}\section*{3. 概形的 Grothendieck 群的补充\footnote{本节结果在\hyperlink{sga6.e05.expose}{第 V} 至 VIII 讲中不会用到.}}

\sgaStableTarget{sga6.e04.prop.3-1-0}\textbf{命题 3.1.0.} 设
\[
\begin{tikzcd}
X' \arrow[r,"g"]\arrow[d,"f'"'] & X\arrow[d,"f"]\\
Y' \arrow[r,"h"'] & Y
\end{tikzcd}
\]
为概形的笛卡尔方块, 其中 \(Y\) 拟紧, \(f\) 拟紧且拟分离. 设 \(E\in\Ob D^b(X)\) 为上同调拟凝聚的复形. 假设 \(X\) 与 \(Y'\) 在 \(Y\) 上 Tor 无关（\hyperlink{sga6.e03.def.1-5}{III 1.5}）. 再假设 \(h\) 有有限 Tor 维数, 或 \(E\) 相对于 \(f\) 有有限平坦振幅. 在这些条件下, 底变换箭头（参见 SGA 4 XVII 4）
\[
 Lh^*\,Rf_*E\longrightarrow Rf'_*\,Lg^*E
\]
有定义, 且为同构.

\emph{证明.} 为定义此箭头, 使用 \(h\) 的平凡对偶公式（SGA 4 XVII 2.3）, 即对 \(L\in\Ob D^-(Y)\) 和 \(M\in\Ob D^+(Y')\), 有典范的函子性同构
\begin{equation}\tag{*}
 \Hom(Lh^*L,M)\xrightarrow{\sim}\Hom(L,Rh_*M)
\end{equation}
利用（\hyperlink{sga6.e03.cor.3-5-2}{III 3.5.2}）和（\hyperlink{sga6.e03.cor.3-7-1}{III 3.7.1}）, 容易验证 \(Lg^*E\in\Ob D^b(X')\) 及 \(Rf_*E\in\Ob D^-(Y)\). 因而在上述公式中取 \(L=Rf_*E\) 和 \(M=Rf'_*Lg^*E\), 因而应用此公式, 问题归结为定义一个箭头
\[
 Rf_*E\longrightarrow Rf_*Rg_*Lg^*E.
\]
取该箭头为如下伴随映射在 \(Rf_*\) 下的像:
\[
 E\longrightarrow Rg_*Lg^*E,
\]
此映射来自 \(g\) 的平凡对偶公式. 至此, 底变换箭头已有定义.

还需证明它为同构. 为此, 首先可假设 \(Y\) 与 \(Y'\) 为仿射概形. 再由 \(X\) 的开覆盖的 Leray 谱序列, 归结为 \(X\) 与 \(X'\) 为仿射概形的情形. 因而有环的余笛卡尔方块
\[
\begin{tikzcd}
B\arrow[r] & B'\\
A\arrow[u]\arrow[r] & A'\arrow[u],
\end{tikzcd}
\]
且可假设 \(E\) 由一个上有界的 \(B\)-模复形定义. \(A'\) 与 \(B\) 在 \(A\) 上的 Tor 无关性写成
\[
 B\otimes_A^{\mathbf L}A'\xrightarrow{\sim}B',
\]
它在 \(D(A')\) 中蕴含
\[
 E\otimes_A^{\mathbf L}A'\xrightarrow{\sim}E\otimes_B^{\mathbf L}B',
\]
从而证明了命题.

将 \hyperlink{sga6.e04.prop.3-1-0}{(3.1.0)} 改用 Grothendieck 群表述, 得到:

\sgaStableTarget{sga6.e04.prop.3-1-1}\textbf{命题 3.1.1.} 在概形的笛卡尔方块
\[
\begin{tikzcd}
X' \arrow[r,"g"]\arrow[d,"f'"'] & X\arrow[d,"f"]\\
Y' \arrow[r,"h"'] & Y,
\end{tikzcd}
\]
中, 假设 \(Y\) 与 \(Y'\) 拟紧, \(X\) 与 \(Y'\) 在 \(Y\) 上 Tor 无关（\hyperlink{sga6.e03.def.1-5}{III 1.5}）, \(f\) 固有且伪凝聚, 并且有限性猜想（\hyperlink{sga6.e03.conj.2-1}{III 2.1}）对 \(f\) 和 \(f'\) 成立; 例如, \(Y\) 与 \(Y'\) 均为 Noether 概形, 或 \(f\) 为射影态射时即是如此. 那么, 若 \(f\) 完美（分别若 \(h\) 有有限 Tor 维数）, 则对 \(x\in K^{\bullet}(X)\)（分别对 \(x\in K_{\bullet}(X)\)）, 有
\[
 h^* f_*(x)=f'_* g^*(x),
\]
其中 \(h^*,g^*\) 由 \hyperlink{sga6.e04.eq.2-7-1}{(2.7.1)}（分别由 \hyperlink{sga6.e04.eq.2-12-1}{(2.12.1)}）定义, \(f_*,f'_*\) 由 \hyperlink{sga6.e04.eq.2-12-3}{(2.12.3)}（分别由 \hyperlink{sga6.e04.eq.2-11-1-1}{(2.11.1.1)}）定义. 注意, 由（\hyperlink{sga6.e03.cor.4-7-1}{III 4.7.1}）, 若 \(f\) 完美, 则 \(f'\) 完美; 由（\hyperlink{sga6.e03.cor.3-5-1}{III 3.5.1}）, 若 \(h\) 有有限 Tor 维数, 则 \(g\) 亦然.

\sgaStableTarget{sga6.e04.para.3-2}\textbf{3.2. \(K^{\bullet}_{\naif}\) 与 \(K_{\bullet}\) 的取极限.} 设 \((X_i,f_{ij})_{i\in I}\) 为拟紧拟分离概形的滤过射影系, 其转移态射为仿射态射（参见 EGA IV 8.2.2）. 置
\[
 X=\varprojlim_i X_i,
\]
用 \(f_i:X\to X_i\) 表示典范态射.

\sgaStableTarget{sga6.e04.para.3-2-1}\textbf{3.2.1.} 同态 \hyperlink{sga6.e04.eq.2-7-2}{(2.7.2)}
\[
 f_i^*:K^{\bullet}(X_i)_{\naif}\longrightarrow K^{\bullet}(X)_{\naif}
\]
由传递公式 \hyperlink{sga6.e04.eq.2-7-3}{(2.7.3)}, 定义同态
\sgaStableTarget{sga6.e04.eq.3-2-1-1}\begin{equation}\tag{3.2.1.1}
 \varinjlim_i K^{\bullet}(X_i)_{\naif}\longrightarrow K^{\bullet}(X)_{\naif}.
\end{equation}
我断言, 此同态为同构. 事实上, 设 \(E\) 为有限型局部自由 \(\mathcal O_X\)-模. 由 EGA IV 8.5.2 和 8.5.5, 存在指标 \(i\) 及有限型局部自由 \(\mathcal O_{X_i}\)-模 \(E_i\), 使得
\[
 E\simeq f_i^*(E_i).
\]
由 EGA IV 8.5.2.5, \(\cl(E_i)\) 在 \(\varinjlim_i K^{\bullet}(X_i)_{\naif}\) 中的像与 \(E_i\) 的选择无关; 将它记为 \(u(E)\). 由 EGA IV 8.5.2 和 8.5.6, 如此定义的映射
\[
 u:\Ob\Loclib(X)\longrightarrow \varinjlim_i K^{\bullet}(X_i)_{\naif}
\]
在 \hyperlink{sga6.e04.para.1-5}{(1.5)} 的意义下可加. 因而它定义同态
\[
 K^{\bullet}(X)_{\naif}\longrightarrow \varinjlim_i K^{\bullet}(X_i)_{\naif},
\]
利用 EGA IV 8.5.2 立即验证, 它为 (3.2.1.1) 的逆. 断言得证.

\sgaStableTarget{sga6.e04.para.3-2-2}\textbf{3.2.2.} 同样可定义典范同态
\sgaStableTarget{sga6.e04.eq.3-2-2-1}\begin{equation}\tag{3.2.2.1}
 \varinjlim_i K^{\bullet}(X_i)\longrightarrow K^{\bullet}(X).
\end{equation}
利用 Deligne 的局部化技巧（SGA 4 XVII）, 它蕴含 \(X\) 上完美复形的范畴等价于 \(X_i\) 上相应完美复形范畴的范畴
\[
 \varinjlim_i \Parf(X_i)
\]
可再次证明 \hyperlink{sga6.e04.eq.3-2-2-1}{(3.2.2.1)} 为同构.

\sgaStableTarget{sga6.e04.para.3-2-3}\textbf{3.2.3.} 假设各 \(X_i\) 及 \(X\) 为 Noether 概形, 且态射 \(f_i\) 平坦（此时各 \(f_{ij}\) 也平坦）. 则同态 \hyperlink{sga6.e04.eq.2-12-1}{(2.12.1)}
\[
 f_i^*:K_{\bullet}(X_i)\longrightarrow K_{\bullet}(X)
\]
定义同构
\sgaStableTarget{sga6.e04.eq.3-2-3-1}\begin{equation}\tag{3.2.2.1}
 \varinjlim_i K_{\bullet}(X_i)\xrightarrow{\sim}K_{\bullet}(X).
\end{equation}
证明与 \hyperlink{sga6.e04.para.3-2-1}{(3.2.1)} 相似, 利用 \hyperlink{sga6.e04.para.2-4}{(2.4)} 中的事实: \(K_{\bullet}(X)\simeq K_{\bullet}(X)_{\naif}\) 可由 \(X\) 上的凝聚层计算.

此外, 由 \hyperlink{sga6.e04.eq.2-7-1}{(2.7.1)} 和 \hyperlink{sga6.e04.eq.2-12-2}{(2.12.2)},\footnote{译校注: 原书第 293 页将此处的 (2.7.1) 印作 (2.7.a).} 各 \(K_{\bullet}(X_i)\) 在 \(K^{\bullet}(X_i)_{\naif}\) 上的模结构与逆像 \(f_{ij}^*\) 相容. 取极限后得到, \(\varinjlim_i K_{\bullet}(X_i)\) 是 \(\varinjlim_i K^{\bullet}(X_i)_{\naif}\) 上的模; 经同构 \hyperlink{sga6.e04.eq.3-2-1-1}{(3.2.1.1)} 和 \hyperlink{sga6.e04.eq.3-2-3-1}{(3.2.2.1)}, 此结构与 \(K_{\bullet}(X)\) 在 \(K^{\bullet}(X)_{\naif}\) 上的模结构相容. 换言之, 有交换图, 其中竖直箭头由 \hyperlink{sga6.e04.eq.3-2-1-1}{(3.2.1.1)} 与 \hyperlink{sga6.e04.eq.3-2-3-1}{(3.2.2.1)} 定义:
\[
\begin{tikzcd}[column sep=tiny]
\varinjlim K^{\bullet}(X_i)_{\naif}\otimes \varinjlim K_{\bullet}(X_i)
\arrow[r]\arrow[d,"\sim"']
&\varinjlim K_{\bullet}(X_i)\arrow[d,"\sim"]\\
K^{\bullet}(X)_{\naif}\otimes K_{\bullet}(X)\arrow[r]
&K_{\bullet}(X).
\end{tikzcd}
\]

\emph{译校注.} 原书在 3.2.2 与 3.2.3 中重复使用式号 (3.2.2.1). 本段所引的是上面的 \(K_{\bullet}\) 同构; 若依段落区分, 其式号应为 (3.2.3.1).

\sgaStableTarget{sga6.e04.para.3-3}\textbf{3.3. 相对 Grothendieck 群.}

\sgaStableTarget{sga6.e04.para.3-3-1}\textbf{3.3.1.} 设
\[
 f:X\longrightarrow Y
\]
为概形的局部有限型态射. 置
\[
 K^{\bullet}(f)=k(D(f)_{\underline{\parf}}),\qquad
 K_{\bullet}(f)=k(D^b(f)_{\coh}),
\]
其中 \(D^b(f)_{\coh}\) 表示相对于 \(f\) 伪凝聚且上同调有界的复形范畴（\hyperlink{sga6.e03.def.1-2}{III 1.2}）, \(D(f)_{\underline{\parf}}\) 表示相对于 \(f\) 有限完美振幅的复形范畴（\hyperlink{sga6.e03.notation.4-2}{III 4.2}）. 在不致混淆时, 有时将 \(K^{\bullet}(f)\)（分别将 \(K_{\bullet}(f)\)）写作 \(K^{\bullet}(X/Y)\)（分别写作 \(K_{\bullet}(X/Y)\)）.

群 \(K^{\bullet}(f)\) 与 \(K_{\bullet}(f)\) 没有其绝对情形 \(K^{\bullet}(X)\) 和 \(K_{\bullet}(X)\) 那样丰富的结构, 因为张量积一般不定义同态
\[
 K^{\bullet}(f)\otimes K^{\bullet}(f)\longrightarrow K^{\bullet}(f),
 \qquad
 K^{\bullet}(f)\otimes K_{\bullet}(f)\longrightarrow K_{\bullet}(f).
\]
另一方面, 它们满足有趣的变性性质, 下面予以回顾.

\sgaStableTarget{sga6.e04.para.3-3-2}\textbf{3.3.2.} 设
\[
\begin{tikzcd}
X\arrow[rr,"f"]\arrow[dr,"h"'] &&Y\arrow[dl,"g"]\\
&S&
\end{tikzcd}
\]
为概形的局部有限型态射的交换三角形.
\begin{enumerate}[label=\textbf{3.3.2.\arabic*.}, leftmargin=4em]
\item 若 \(g\) 伪凝聚, 则有
\[
 K_{\bullet}(f)=K_{\bullet}(h)
\]
（\hyperlink{sga6.e03.cor.1-12}{III 1.12}）. 若 \(g\) 光滑且 \(X\) 拟紧, 则有
\[
 K^{\bullet}(f)=K^{\bullet}(h)
\]
（\hyperlink{sga6.e03.prop.4-4}{III 4.4.1}）.

\item 由（\hyperlink{sga6.e03.prop.1-11}{III 1.11} 及 \hyperlink{sga6.e03.prop.4-5}{4.5}）, 函子 \(\otimes_Y^{\mathbf L}\) 定义配对
\[
 D(f)_{\underline{\parf}}\times D^b(g)_{\coh}\longrightarrow D^b(h)_{\coh},
\]
\[
 D(f)_{\underline{\parf}}\times D(g)_{\underline{\parf}}\longrightarrow D(h)_{\underline{\parf}},
\]
它们诱导同态
\[
 K^{\bullet}(f)\otimes K_{\bullet}(g)\longrightarrow K_{\bullet}(h),
\]
\[
 K^{\bullet}(f)\otimes K^{\bullet}(g)\longrightarrow K^{\bullet}(h).
\]
特别地, 取 \(X=Y\) 和 \(f=\id\), 得到 \(K_{\bullet}(g)\) 是 \(K^{\bullet}(Y)\)-模. 另一方面, 若 \(f\) 完美, 则乘以
\[
 \cl(\mathcal O_X)\in K^{\bullet}(f)
\]
定义 \(K^{\bullet}(Y)\)-线性同态
\[
 f^*:K_{\bullet}(Y/S)\longrightarrow K_{\bullet}(X/S),
\]
以及
\[
 f^*:K^{\bullet}(Y/S)\longrightarrow K^{\bullet}(X/S).
\]

\item 假设 \(S\) 拟紧, \(X\) 与 \(Y\) 在 \(S\) 上为有限型, 且 \(f\) 固有. 则只要有限性猜想（\hyperlink{sga6.e03.conj.2-1}{III 2.1}）成立, 特别地, 若 \(Y\) 为 Noether 概形或 \(f\) 为射影态射, 函子 \(Rf_*\) 将 \(D^b(h)_{\coh}\) 映入 \(D^b(g)_{\coh}\), 因而定义同态
\[
 f_*:K_{\bullet}(X/S)\longrightarrow K_{\bullet}(Y/S).
\]
若进一步 \(f\) 完美, 则 \(Rf_*\) 将 \(D(h)_{\underline{\parf}}\) 映入 \(D(g)_{\underline{\parf}}\)（\hyperlink{sga6.e03.prop.4-8}{III 4.8}）, 并定义同态
\[
 f_*:K^{\bullet}(X/S)\longrightarrow K^{\bullet}(Y/S).
\]
由投影公式（\hyperlink{sga6.e03.prop.3-7}{III 3.7}）得到
\[
 f_*(a x)=a f_*(x),
\]
其中 \(f\) 固有, \(a\in K^{\bullet}(S)\), \(x\in K_{\bullet}(X)\); 分别地, \(f\) 固有且完美, \(a\in K_{\bullet}(S)\), \(x\in K^{\bullet}(X/S)\).

\item 假设 \(S\) 为 Noether 概形, \(X\) 与 \(Y\) 在 \(S\) 上为有限型, 且 \(f\) 可紧化并且完美. 则由（\hyperlink{sga6.e03.cor.4-9-1}{III 4.9.1}）, 函子 \(f^!\) 将 \(D(g)_{\underline{\parf}}\) 映入 \(D(h)_{\underline{\parf}}\), 并定义同态
\[
 f^!:K^{\bullet}(Y/S)\longrightarrow K^{\bullet}(X/S).
\]
\end{enumerate}

\sgaStableTarget{sga6.e04.para.3-3-3}\textbf{3.3.3.} 在（\hyperlink{sga6.e03.prop.4-7}{III 4.7}）的假设下, 函子 \(\otimes_S^{\mathbf L}\) 定义配对
\[
 D(f_1)_{\underline{\parf}}\times D^b(f_2)_{\coh}\longrightarrow D^b(f)_{\coh},
\]
\[
 D(f_1)_{\underline{\parf}}\times D(f_2)_{\underline{\parf}}\longrightarrow D(f)_{\underline{\parf}},
\]
它们诱导同态
\[
 K^{\bullet}(f_1)\otimes K_{\bullet}(f_2)\longrightarrow K_{\bullet}(f),
\]
\[
 K^{\bullet}(f_1)\otimes K^{\bullet}(f_2)\longrightarrow K^{\bullet}(f).
\]
当然, 当 \(X_i=Y_i=S\) 时, 就得到 \(K_{\bullet}(S)\) 上的 \(K^{\bullet}(S)\)-模结构.

若 \(X\) 在 \(S\) 上平坦, 则对 \(n\geq1\) 得到同态
\[
 K^{\bullet}(X/S)^{\otimes n}\longrightarrow K^{\bullet}\bigl((X/S)^n\bigr),
\]
它们在某种意义上应当代替 \(K^{\bullet}(X/S)\) 所缺少的环结构.

\sgaStableTarget{sga6.e04.bibliography}\subsection*{参考文献}
\begin{enumerate}[label={[\arabic*]}]
\item Deligne, P., papiers secrets（未公开手稿）.
\item Grauert, H., \emph{Ein Theorem der Analytischen Garbentheorie}, Publ. Math. IHES, no. 5 (1960).
\item Grothendieck, A., \emph{Théorèmes de finitude pour la cohomologie des faisceaux}, Bull. Soc. Math. France, 84, 1956, pp. 1--7.
\item[{[V]}] Verdier, J.-L., \emph{Catégories dérivées}, IHES (1964), 油印本.
\end{enumerate}



% 第 V 讲所需宏
\providecommand{\Chat}{\widehat C}
\providecommand{\Ghat}{\widehat G}
\providecommand{\Gg}{\widehat G_g}
\providecommand{\Ga}{\widehat G_a}
\providecommand{\Grhat}{\widehat{\operatorname{Gr}}}
\providecommand{\Abfil}{\mathrm{Ab\,fil}}
\providecommand{\Kalgfil}{K\text{-}\mathrm{Alg.fil.}}
\providecommand{\Hom}{\operatorname{Hom}}
\providecommand{\Sym}{\mathfrak S}
\providecommand{\id}{\operatorname{Id}}
\providecommand{\Log}{\operatorname{Log}}
\providecommand{\Toddc}{\operatorname{Todd}}
\providecommand{\Ob}{\operatorname{Ob}}
\providecommand{\cG}{\mathcal G}
\providecommand{\cCzero}{\mathcal C_0}
\providecommand{\Fil}{\operatorname{Fil}}


